\documentclass[../../../main.tex]{subfiles}

\begin{document}

\chapter{Multivariable Functions and Continuity}

Finally, we are going to discuss about multivariable calculus!
Before we get too excited, we need to discuss some of the fundamental
topics on multivariable functions and their limits, continuity, and
differentiability. In this note, we are going to discuss fundamental
and necessary topics before investigating more into other calculus
topics.

\section{Multivariable Functions and Sets}

As always, let's get started with definitions and common notations.
Consider the following notation.
\[
\mathbb{R}^n
\coloneqq \{ (x_1, x_2, \ldots, x_n) \mid x_i \in \mathbb{R} \}
\]
As you can notice from the equation, $\mathbb{R}^n$ represents an
$n$-tuple composed of real numbers. Some examples include
$\mathbb{R}^2$ and $\mathbb{R}^3$, which represents the coordinates in
2-space and 3-space respectively.

\begin{important}
Given this notation, we could naturally define a multivariable
function $f(x_1, x_2, \ldots, x_n)$ as a function that maps
the elements of $D \subset \mathbb{R}^n$ to a single real number.
Here, the domain $D$ is represented as $\dom(f)$ and the range
is denoted as $\ran(f)$.
\end{important}

If the domain is not explicitly stated for a function, we could just
assume it to be the largest subset of $\mathbb{R}^n$ such that the
function $f$ is defined.

An interesting question to consider is how do you graph functions like
$f(x, y)$. The answer to this question is pretty straightforward given
that the graph of $f(x)$ is defined as $\{ (x, y) \mid y = f(x) \}$.
For $f(x, y)$, we could simply extend it as $\{ (x, y, z) \mid
z = f(x, y) \}$. For functions with more variables, it is impossible,
or very difficult, to visually graph them, but there is a strategy
to get a glimpse on their appearance. To get to the strategy,
consider the following definitions.

\begin{definition}{}{}
Given a function $f(x, y)$, the \textbf{level curve} of $f$ of height
$k$ is defined as the set $\{ (x, y) \mid f(x, y) = k \}$.
The collection of different level curve of $f$ is known as a
\textbf{contour plot}, or \textbf{contour map} of $f$.
\end{definition}

Recall that when we wanted to see how paraboloids look like,
we drew many traces of the surface. This is exactly what it is!
With the same idea, we can get a glimpse of how a complex functions
in 3-space look like. Similarly, we could do the same for functions
in $4$-space. Although it is almost impossible and very difficult,
we can use this method to approximate its appearance.

Continuing from the discussion on multivariable functions, let's now
turn our attention to their domains.



\end{document}


